\textbf{Phần 1. Trắc nghiệm lựa chọn}
\setcounterpageref{ex}{0}

\begin{ex} %[3X1V75]
Để ghi chép số lượng xe đi qua cổng trường trong 15 phút, bạn An dùng các vạch kiểm đếm (mỗi nhóm gồm 4 vạch dọc và 1 vạch chéo ngang biểu thị 5 xe). An đếm được 6 nhóm vạch và 3 vạch rời. Số xe An đếm được là
\choice
{\True $33$ xe}
{$30$ xe}
{$35$ xe}
{$28$ xe}
\loigiai{
Số xe là: $5 \times 6 + 3 = 33$ xe.
\textbf{Đáp án của câu hỏi là A}
}
\end{ex}

\begin{ex} %[3X1V75]
Bảng thống kê thời gian tự học ở nhà mỗi ngày của 4 bạn: An: 60 phút; Bình: 45 phút; Cúc: 90 phút; Dũng: 75 phút. Bạn dành nhiều thời gian tự học nhất là
\choice
{An}
{Dũng}
{\True Cúc}
{Bình}
\loigiai{
Cúc học 90 phút, nhiều nhất.
\textbf{Đáp án của câu hỏi là C}
}
\end{ex}

\begin{ex} %[3X1V75]
Thời gian tự học của Cúc nhiều hơn Bình bao nhiêu phút?
\choice
{$30$ phút}
{\True $45$ phút}
{$15$ phút}
{$60$ phút}
\loigiai{
$90 - 45 = 45$ phút.
\textbf{Đáp án của câu hỏi là B}
}
\end{ex}

\begin{ex} %[3X1V75]
Thống kê chiều cao 4 bạn học sinh: Nam cao $132\text{ cm}$, Hoa cao $128\text{ cm}$, Tuấn cao $135\text{ cm}$, Lan cao $130\text{ cm}$. Bạn thấp nhất là
\choice
{Lan}
{Nam}
{Tuấn}
{\True Hoa}
\loigiai{
Hoa cao $128\text{ cm}$, thấp nhất.
\textbf{Đáp án của câu hỏi là D}
}
\end{ex}

\begin{ex} %[3X1V75]
Chênh lệch chiều cao giữa bạn cao nhất và bạn thấp nhất trong nhóm trên là
\choice
{\True $7\text{ cm}$}
{$5\text{ cm}$}
{$6\text{ cm}$}
{$8\text{ cm}$}
\loigiai{
Tuấn cao nhất ($135\text{ cm}$), Hoa thấp nhất ($128\text{ cm}$). Chênh lệch: $135 - 128 = 7\text{ (cm)}$.
\textbf{Đáp án của câu hỏi là A}
}
\end{ex}

\begin{ex} %[3X1V75]
Số học sinh tham gia câu lạc bộ bơi của các lớp khối 3: 3A có 14 bạn; 3B có 16 bạn; 3C có 15 bạn; 3D có 15 bạn. Tổng số học sinh tham gia của cả khối là
\choice
{$55$ bạn}
{$65$ bạn}
{\True $60$ bạn}
{$58$ bạn}
\loigiai{
$14 + 16 + 15 + 15 = 60$ bạn.
\textbf{Đáp án của câu hỏi là C}
}
\end{ex}

\begin{ex} %[3X1V75]
Một bạn gieo đồng xu 20 lần và ghi lại: có 12 lần xuất hiện mặt Sấp. Số lần xuất hiện mặt Ngửa là
\choice
{$10$ lần}
{\True $8$ lần}
{$12$ lần}
{$6$ lần}
\loigiai{
Số lần mặt Ngửa: $20 - 12 = 8$ lần.
\textbf{Đáp án của câu hỏi là B}
}
\end{ex}

\begin{ex} %[3X1V75]
Ghi chép số quả trứng gà thu hoạch trong 5 ngày: Ngày 1: 45 quả; Ngày 2: 50 quả; Ngày 3: 48 quả; Ngày 4: 52 quả; Ngày 5: 55 quả. Ngày thu hoạch được nhiều trứng nhất là
\choice
{Ngày 4}
{Ngày 2}
{Ngày 3}
{\True Ngày 5}
\loigiai{
Ngày 5 thu hoạch được 55 quả, nhiều nhất.
\textbf{Đáp án của câu hỏi là D}
}
\end{ex}

\begin{ex} %[3X1V75]
Tổng số trứng gà thu hoạch được trong 5 ngày đó là
\choice
{\True $250$ quả}
{$240$ quả}
{$260$ quả}
{$245$ quả}
\loigiai{
Tổng số trứng: $45 + 50 + 48 + 52 + 55 = 250$ quả.
\textbf{Đáp án của câu hỏi là A}
}
\end{ex}

\begin{ex} %[3X1V75]
Khi thực hiện một cuộc phỏng vấn nhanh về môn thể thao yêu thích, công cụ ghi chép nhanh và hiệu quả nhất tại chỗ là
\choice
{Thước kẻ}
{Vẽ tranh}
{\True Bảng kiểm đếm bằng các vạch gạch}
{Máy tính để bàn cồng kềnh}
\loigiai{
Ghi chép bằng vạch kiểm đếm là phương pháp chuẩn xác và nhanh nhất.
\textbf{Đáp án của câu hỏi là C}
}
\end{ex}

\begin{ex} %[3X1V75]
Kết quả bình bầu cán bộ lớp: Mai được 28 phiếu, Lan được 24 phiếu, Cường được 20 phiếu. Số phiếu của Mai nhiều hơn Cường bao nhiêu phiếu?
\choice
{$4$ phiếu}
{\True $8$ phiếu}
{$6$ phiếu}
{$10$ phiếu}
\loigiai{
$28 - 20 = 8$ phiếu.
\textbf{Đáp án của câu hỏi là B}
}
\end{ex}

\begin{ex} %[3X1V75]
Thống kê số lượng cây hoa trồng trong bồn hoa lớp: 15 cây hoa hồng, 20 cây hoa cúc, 10 cây hoa đồng tiền. Tổng số cây hoa đã trồng là
\choice
{$40$ cây}
{$50$ cây}
{$35$ cây}
{\True $45$ cây}
\loigiai{
$15 + 20 + 10 = 45$ cây.
\textbf{Đáp án của câu hỏi là D}
}
\end{ex}

\textbf{Phần 2. Trắc nghiệm đúng sai}
\setcounterpageref{ex}{0}

\begin{ex} %[3X1V75]
Xét tính đúng sai của việc thực hành trải nghiệm thống kê:
\choiceTF
{\True Kiểm đếm bằng vạch giúp đếm số lượng nhanh và tránh nhầm lẫn.}
{Chỉ có thể thu thập số liệu bằng cách cân đo, không thể đếm.}
{\True Bảng số liệu thống kê phản ánh trung thực kết quả thu thập.}
{Kết quả thực nghiệm xác suất sau 10 lần gieo đồng xu luôn luôn có 5 lần sấp và 5 lần ngửa.}
\loigiai{
\begin{itemize}
    \item Ý a) Đúng: Lợi ích kiểm đếm.
    \item Ý b) Sai: Sai vì đếm số lượng là phương pháp thống kê cơ bản.
    \item Ý c) Đúng: Yêu cầu khoa học.
    \item Ý d) Sai: Sai vì kết quả ngẫu nhiên có thể dao động.
\end{itemize}
\begin{itemize}
    \item a) \textbf{Đúng}.
    \item b) \textbf{Sai}.
    \item c) \textbf{Đúng}.
    \item d) \textbf{Sai}.
\end{itemize}
}
\end{ex}

\begin{ex} %[3X1V75]
Bảng thống kê số tiết đọc sách thư viện của 3 bạn trong tháng: An đọc 12 tiết, Bình đọc 16 tiết, Cúc đọc 20 tiết:
\choiceTF
{\True Cúc đọc nhiều tiết sách nhất trong ba bạn.}
{\True Cả ba bạn đọc tất cả 48 tiết sách.}
{Bình đọc gấp đôi số tiết của An.}
{\True Số tiết đọc của Cúc nhiều hơn An là 8 tiết.}
\loigiai{
\begin{itemize}
    \item Ý a) Đúng: 20 tiết nhiều nhất.
    \item Ý b) Đúng: $12 + 16 + 20 = 48$ tiết.
    \item Ý c) Sai: Sai vì $12 \times 2 = 24 \neq 16$.
    \item Ý d) Đúng: $20 - 12 = 8$ tiết.
\end{itemize}
\begin{itemize}
    \item a) \textbf{Đúng}.
    \item b) \textbf{Đúng}.
    \item c) \textbf{Sai}.
    \item d) \textbf{Đúng}.
\end{itemize}
}
\end{ex}

\textbf{Phần 3. Trả lời ngắn}
\setcounterpageref{ex}{0}

\begin{ex} %[3X1V75]
Kiểm đếm số học sinh thích màu đỏ: có 7 nhóm vạch (mỗi nhóm 5 vạch) và 2 vạch rời. Tính tổng số học sinh thích màu đỏ.
\shortans{37}
\loigiai{
\textbf{Đáp số:} $37$\\
$5 \times 7 + 2 = 37$ bạn.
}
\end{ex}

\begin{ex} %[3X1V75]
Khối 3 có 4 lớp, mỗi lớp quyên góp được 25 quyển truyện tranh. Cả khối 3 quyên góp được bao nhiêu quyển truyện tranh?
\shortans{100}
\loigiai{
\textbf{Đáp số:} $100$\\
$25 \times 4 = 100$ quyển.
}
\end{ex}

\begin{ex} %[3X1V75]
Thực nghiệm gieo xúc xắc 30 lần, xuất hiện mặt 6 chấm 8 lần. Hỏi các mặt còn lại xuất hiện bao nhiêu lần?
\shortans{22}
\loigiai{
\textbf{Đáp số:} $22$\\
$30 - 8 = 22$ lần.
}
\end{ex}

\begin{ex} %[3X1V75]
Bạn Lan đo nhiệt độ trong 3 ngày: Ngày 1: $28^\circ\text{C}$; Ngày 2: $30^\circ\text{C}$; Ngày 3: $32^\circ\text{C}$. Hỏi nhiệt độ ngày 3 hơn ngày 1 bao nhiêu độ C?
\shortans{4}
\loigiai{
\textbf{Đáp số:} $4$\\
$32 - 28 = 4^\circ\text{C}$.
}
\end{ex}

\begin{ex} %[3X1V75]
Số điểm thưởng của 3 tổ: Tổ 1 được 40 điểm, Tổ 2 được 50 điểm, Tổ 3 được 60 điểm. Tính tổng số điểm thưởng của cả 3 tổ.
\shortans{150}
\loigiai{
\textbf{Đáp số:} $150$\\
$40 + 50 + 60 = 150$ điểm.
}
\end{ex}

\begin{ex} %[3X1V75]
Trong lớp có 36 học sinh, có 24 bạn tham gia học bơi. Hỏi có bao nhiêu bạn chưa tham gia học bơi?
\shortans{12}
\loigiai{
\textbf{Đáp số:} $12$\\
$36 - 24 = 12$ bạn.
}
\end{ex}
